\chapter{起步：临床研究的结构与运作}
\label{ch:01}
\chapterauthors{Warren S. Browner、Thomas B. Newman、Mark J. Pletcher}
\tikzset{dcrhead/.style={draw=blue!50!cyan,fill=blue!48!cyan!65,text=white,
 font=\small\bfseries\heiti,minimum height=9mm,align=center},
 dcrbody/.style={draw=blue!45!cyan,text=dcrblue,align=left,font=\footnotesize,
 inner sep=2mm},dcrresult/.style={draw=red!60!black,thick,align=center,
 font=\small\bfseries\heiti,minimum height=13mm},
 dcrarr/.style={-{Latex[length=3mm,width=2.5mm]},line width=1.6pt},
 dcrvalid/.style={text=black!55,font=\small\bfseries\heiti,align=center}}
\newcommand{\chonebasicflow}{%
\begin{tikzpicture}[x=1mm,y=1mm]
\foreach \x/\name/\txt in {0/q/研究问题,48/p/研究计划,96/a/实际研究}
 \node[dcrhead,text width=32mm] (\name) at (\x,0) {\txt};
\foreach \x/\name/\txt in {0/u/当前认识,48/c/研究结论,96/r/研究结果}
 \node[dcrresult,text width=32mm] (\name) at (\x,-39) {\txt};
\draw[dcrarr,blue!60] (q)--node[above,font=\small]{设计}(p);
\draw[dcrarr,blue!60] (p)--node[above,font=\small]{实施}(a);
\draw[dcrarr,red!65!black] (r)--node[above,font=\small]{解释}(c);
\draw[dcrarr,red!65!black] (c)--node[above,font=\small]{推断}(u);
\draw[dcrarr,blue!65!red] (a.east)--++(7,0)|-node[pos=.25,right,rotate=90,font=\small]{分析}(r.east);
\draw[dcrarr,red!65!blue] (u.west)--++(-7,0)|-node[pos=.25,left,rotate=90,font=\small]{思考}(q.west);
\end{tikzpicture}}
\newcommand{\chonedesignflow}{%
\begin{tikzpicture}[x=1mm,y=1mm]
\node[dcrhead,text width=49mm] (q) at (0,0) {研究问题};
\node[dcrhead,text width=49mm] (p) at (84,0) {研究计划};
\node[dcrbody,text width=49mm,minimum height=58mm,anchor=north] (qb) at (q.south)
 {\textbf{目标人群}\\全体 70 岁及以上、认知功能正常的美国人\\[5mm]
 \textbf{目标现象}\\过去一个月每天摄入咖啡因的情况};
\node[dcrbody,text width=49mm,minimum height=58mm,anchor=north] (pb) at (p.south)
 {\textbf{预定样本}\\70 岁及以上、无痴呆或认知功能障碍诊断的医保计划成员\\[5mm]
 \textbf{预定变量}\\自报过去一个月的咖啡、茶、含咖啡因饮料及高能量饮料摄入情况};
\node[dcrresult,text width=49mm] (u) at (0,-102) {当前认识};
\node[dcrresult,text width=49mm] (c) at (84,-102) {研究结论};
\draw[dcrarr,blue!65] (qb.east)--node[above,font=\small]{设计}(pb.west);
\draw[dcrarr,red!65!black] (c)--node[above,font=\small]{推断}(u);
\node[dcrvalid] (v) at (42,-78) {外部效度};
\draw[dcrarr,black!45] (v.north)--(42,-35);
\draw[dcrarr,black!45] (v.south)--(42,-102);
\node[text=dcrgold,align=center,font=\small\bfseries] (e) at (100,-78) {选择与误差\\影响……};
\draw[dcrarr,black!45] (e.west)--(v.east);
\end{tikzpicture}}
\newcommand{\choneimplementationflow}{%
\begin{tikzpicture}[x=1mm,y=1mm]
\node[dcrhead,text width=48mm] (p) at (0,0) {研究计划};
\node[dcrhead,text width=48mm] (a) at (80,0) {实际研究};
\node[dcrbody,text width=48mm,minimum height=60mm,anchor=north] (pb) at (p.south)
 {\textbf{预定样本}\\70 岁及以上、无痴呆或认知功能障碍诊断的医保计划成员\\[5mm]
 \textbf{预定变量}\\自报过去一个月的咖啡、茶、含咖啡因饮料及高能量饮料摄入情况};
\node[dcrbody,text width=48mm,minimum height=60mm,anchor=north] (ab) at (a.south)
 {\textbf{实际样本}\\70 岁及以上、无痴呆或认知功能障碍诊断且寄回问卷的医保计划成员\\[4mm]
 \textbf{实际测量}\\关于过去一个月咖啡、茶、含咖啡因饮料及高能量饮料摄入情况的回答};
\node[dcrresult,text width=48mm] (c) at (0,-105) {研究结论};
\node[dcrresult,text width=48mm] (r) at (80,-105) {研究结果};
\draw[dcrarr,blue!65] (pb.east)--node[above,font=\small]{实施}(ab.west);
\draw[dcrarr,red!65!black] (r)--node[above,font=\small]{解释}(c);
\draw[dcrarr,blue!65!red] (ab.east)--++(9,0)|-node[pos=.28,right,font=\small]{分析}(r.east);
\node[dcrvalid] (v) at (40,-81) {内部效度};
\draw[dcrarr,black!45] (v.north)--(40,-35);
\draw[dcrarr,black!45] (v.south)--(40,-105);
\draw[dcrarr,black!45] (v.east)--(115,-81);
\node[text=dcrgold,align=center,font=\small\bfseries] (e) at (-9,-81) {选择与误差\\影响……};
\draw[dcrarr,black!45] (e.east)--(v.west);
\end{tikzpicture}}
\newcommand{\chonefullflow}{%
\begin{tikzpicture}[x=1mm,y=1mm]
\node[dcrhead,text width=32mm] (q) at (0,0) {研究问题};
\node[dcrhead,text width=32mm] (p) at (49,0) {研究计划};
\node[dcrhead,text width=32mm] (a) at (98,0) {实际研究};
\node[dcrbody,text width=32mm,minimum height=80mm,anchor=north] (qb) at (q.south)
 {\textbf{目标人群}\\全体 70 岁及以上、认知功能正常的美国人\\[5mm]
 \textbf{目标现象}\\过去一个月每天摄入咖啡因的情况};
\node[dcrbody,text width=32mm,minimum height=80mm,anchor=north] (pb) at (p.south)
 {\textbf{预定样本}\\70 岁及以上、无痴呆或认知功能障碍诊断的医保计划成员\\[4mm]
 \textbf{预定变量}\\自报过去一个月的咖啡、茶、含咖啡因饮料及高能量饮料摄入情况};
\node[dcrbody,text width=32mm,minimum height=80mm,anchor=north] (ab) at (a.south)
 {\textbf{实际样本}\\70 岁及以上、无痴呆或认知功能障碍诊断且寄回问卷的医保计划成员\\[3mm]
 \textbf{实际测量}\\关于过去一个月咖啡、茶、含咖啡因饮料及高能量饮料摄入情况的回答};
\node[dcrresult,text width=32mm] (u) at (0,-135) {当前认识};
\node[dcrresult,text width=32mm] (c) at (49,-135) {研究结论};
\node[dcrresult,text width=32mm] (r) at (98,-135) {研究结果};
\draw[dcrarr,blue!65] (qb.east)--node[above,font=\footnotesize]{设计}(pb.west);
\draw[dcrarr,blue!65] (pb.east)--node[above,font=\footnotesize]{实施}(ab.west);
\draw[dcrarr,red!65!black] (r)--node[above,font=\footnotesize]{解释}(c);
\draw[dcrarr,red!65!black] (c)--node[above,font=\footnotesize]{推断}(u);
\draw[dcrarr,blue!65!red] (ab.east)--++(6,0)|-node[pos=.3,right,rotate=90,font=\footnotesize]{分析}(r.east);
\draw[dcrarr,red!65!blue] (u.west)--++(-6,0)|-node[pos=.3,left,rotate=90,font=\footnotesize]{思考}(qb.west);
\node[dcrvalid] (ev) at (24.5,-108) {外部效度};
\node[dcrvalid] (iv) at (73.5,-108) {内部效度};
\node[text=dcrgold,align=center,font=\footnotesize\bfseries] (e) at (49,-108) {选择与误差\\影响……};
\draw[dcrarr,black!45] (e.west)--(ev.east);
\draw[dcrarr,black!45] (e.east)--(iv.west);
\draw[dcrarr,black!45] (ev.north)--(24.5,-35);
\draw[dcrarr,black!45] (ev.south)--(u.north east);
\draw[dcrarr,black!45] (iv.north)--(73.5,-35);
\draw[dcrarr,black!45] (iv.south)--(c.north east);
\draw[dcrarr,black!45] (iv.east)--(121,-108);
\end{tikzpicture}}

\sourcepage{16}
本章从两个角度介绍临床研究（clinical research），并展开贯穿全书的两条主线。第一条是研究的“解剖结构”（anatomy），即研究计划（study plan）中的具体要素，例如研究问题、设计、研究参与者（study participants）、测量和样本量计算等。研究者的目标，是以使项目可行、高效并提高研究发现效度（validity）的方式，设计这些组成部分。

另一条主线是研究的“生理功能”（physiology），也就是研究如何运作。一项研究是否有用，取决于它能否对研究中发生的情况作出有效结论，以及这些结论能否进一步推广到更广阔的世界。目标是尽量减少可能威胁这两个过程的误差。

把这两条主线分开是人为的，正如不理解人体生理，就难以理解人体解剖一样。不过，这样区分有一个好处：让我们更容易思考复杂的问题。

\section{研究的结构：它由什么组成}
研究项目的结构写在\term{研究方案（protocol）}中，也就是研究的书面计划。人们熟知研究方案用于申请科研经费和机构审查委员会（Institutional Review Board，IRB）批准，但它还有至关重要的科学作用：帮助研究者以有逻辑、有重点且高效的方式组织研究。表 1.1 列出研究方案的组成部分。这里先介绍全貌，后续各章逐项展开，到第 20 章再将完成的部分组合起来。

\sourcepage{17}
\begin{center}\small
\captionof*{table}{\textbf{表 1.1\quad 研究的结构：研究计划}}
\begin{tabularx}{\linewidth}{L{62mm}Y}
\toprule\textbf{设计要素}&\textbf{目的}\\\midrule
研究问题&研究要回答哪些问题？\\
背景与意义&为什么这些问题重要？\\
设计：研究类型&研究如何组织？\\
参与者：选择标准、抽样设计&参与者是谁，如何选择？\\
变量：预测变量、混杂变量、结局变量&将进行哪些测量？\\
统计学问题：假设、样本量、分析方法&研究规模多大，如何分析？\\\bottomrule
\end{tabularx}\end{center}

\subsection{研究问题}
\term{研究问题（research question）}\footnote{原书以红色标出的术语，在书末“术语释义”中有定义。}是研究的目标，也就是研究者希望解决的不确定之处。研究问题常从一个宽泛的关切开始，需要进一步收窄为具体、可研究的问题。例如，“咖啡因（caffeine）是否影响认知功能（cognitive function）”是一个不错的起点，但在着手制订研究计划之前，还必须进一步聚焦。这通常意味着把问题拆成更具体的部分，选出其中一两个作为研究方案的中心。可能的研究问题包括：
\begin{itemize}
\item 咖啡因对认知功能的影响是短期的、长期的，还是两者都有？
\item 长期摄入咖啡因是否降低发生痴呆（dementia）的风险？
\item 怎样测量咖啡因摄入量最好？
\sourcepage{18}
\item 摄入咖啡因是否有害？
\item 咖啡因来自咖啡、茶还是含咖啡因的软饮料，是否有影响？
\end{itemize}
关于这一主题，当然还可以提出很多研究问题。研究者的任务，是选择一个可行、重要、新颖且符合伦理的问题。这四项特点有一个便于记忆的缩写——FINE，第 2 章将详细讨论。

\subsection{背景与意义}
要提出有意义的研究问题，必须了解拟研究的领域。与专家交谈并全面查阅文献，常会促使研究者修改问题。因此，设计临床研究项目的第一步，是确定有关主题哪些已经知道、哪些尚不知道；然后思考如何在现有知识之上有所增益。这一步需要长时间深入思考，并与同事和导师讨论重要但尚未回答的问题。

这些想法最终会整合成简明的“背景与意义”部分，说明开展研究的理由。研究者在其中引用相关既往研究，尤其是自己和同事的研究，着重指出那些研究的问题及尚存的不确定之处。同时还要具体说明，拟议研究的发现如何通过产生新的科学知识、改变临床实践，或影响实践指南与卫生政策，帮助解决这些不确定之处。

\subsection{设计}
不存在适用于所有研究问题的最佳设计。研究者应根据所选问题和可用资源，判断哪一种设计最合适。一个基本决定是：在\term{临床试验（clinical trial）}中实施\term{干预（intervention）}并考察其效果，还是在\term{观察性研究（observational study）}中仅对参与者进行测量（表 1.2）。在各种临床试验中，\term{随机盲法试验（randomized blinded trial）}通常最为理想；\sourcepage{19}不过，对某些干预而言，非随机或非盲法设计偶尔可能是唯一可行的选择。

\begingroup\small
\begin{longtable}{L{30mm}L{49mm}L{66mm}}
\caption*{\textbf{表 1.2\quad 探究咖啡因摄入是否降低痴呆风险的临床研究设计示例}}\\
\toprule\textbf{流行病学设计}&\textbf{主要特点}&\textbf{示例}\\\midrule\endfirsthead
\toprule\textbf{流行病学设计}&\textbf{主要特点}&\textbf{示例}\\\midrule\endhead
\multicolumn{3}{c}{\textbf{观察性设计}}\\
横断面研究&在一个时间点检查一组人&研究者访谈一组参与者，了解目前及既往咖啡因摄入情况，并分析这些结果与认知功能测验评分的相关性。\\
队列研究&在研究开始时确定一组人（通常尚无所研究的结局），随后随时间随访&她测量一组认知功能正常者的咖啡因摄入量，并在随访时检查他们，观察摄入咖啡因者是否更少发生痴呆。\\
病例对照研究&依据有无某一结局选择两组人&她比较一组痴呆患者（“病例”）与一组认知功能正常者（“对照”）既往的咖啡因摄入量。\\
\multicolumn{3}{c}{\textbf{临床试验设计}}\\
随机盲法试验&经随机过程产生两组人，干预采用盲法&她将认知功能正常的参与者随机分配至咖啡因补充剂组或外观完全相同的安慰剂组，再对两组随访数年，观察痴呆发生率。\\\bottomrule
\end{longtable}\endgroup
\sourcepage{20}
在观察性研究中，两种常见设计是\term{横断面研究（cross-sectional study）}和\term{队列研究（cohort study）}：前者在一次观察中获得资料，后者对一组参与者随时间进行观察。队列研究又可分为从现在开始、向未来随访的\term{前瞻性队列研究（prospective cohort study）}，以及考察过去收集信息的\term{回顾性队列研究（retrospective cohort study）}，不过这些界限并不总是截然分明。第三种常见的观察性设计是\term{病例对照研究（case-control study）}，研究者将已发生所关注结局的一组人（\term{病例，cases}）与未发生该结局的另一组人（\term{对照，controls}）进行比较。

对于每一个研究问题，都需要判断哪种设计最有效率，能够获得可靠答案。随机盲法试验常被认为是确立因果关系、确定干预效力（efficacy）的最佳设计；但在许多情况下，观察性研究更合适，或是唯一可行的选择。例如，随机分配人们开始或停止喝咖啡，并不容易。病例对照研究成本相对较低，又适合罕见结局，因此对某些问题很有吸引力。这些问题将在第 8—14 章讨论，各章分别介绍一组特定的设计。

一个主题的研究，通常从所谓\term{描述性研究（descriptive studies）}的观察性研究开始。这些研究探索总体状况，例如描述人群中健康相关特征及疾病的分布：
\begin{itemize}
\item 在 70 岁及以上成年人中，摄入咖啡因有多普遍？
\item 老年人中有多大比例认知功能异常？
\end{itemize}
医学检验研究，例如某种新生物标志物的血清水平是否与认知功能障碍相关，是\sourcepage{21}一类特殊的描述性研究（见第 13 章）。

描述性研究之后，或与其同时，通常会开展评价变量间关联的\term{分析性研究（analytic studies）}，通常用以推断因果关系：
\begin{itemize}
\item 每天平均咖啡因摄入量是否与认知功能测验的表现相关？
\item 与认知功能正常的对照相比，痴呆患者既往是否较少经常摄入咖啡因？
\end{itemize}
最后一步常是开展临床试验，确立干预的效果：
\begin{itemize}
\item 随机分配服用咖啡因补充剂的老年志愿者，发生痴呆的风险是否低于被分配服用安慰剂（placebo）者？
\end{itemize}
在围绕特定问题开展的一系列研究中，临床试验通常较晚进行，因为它往往更困难、费用更高，更可能使参与者遭受伤害，而且通常回答的是由观察性研究结果引出的、范围较窄的问题。

用一句话概括设计和研究问题，有助于说明一项研究的性质。如果研究有两个主要阶段，就应分别提到两个阶段的设计。
\begin{studyquote}
这是一项研究 60—74 岁人群咖啡因摄入与认知功能关联的横断面研究，随后以前瞻性队列研究考察咖啡因摄入是否与今后较低的认知下降速度相关。
\end{studyquote}
有些设计不容易归入上述类别，要在一句话中用适当的详略程度描述它们颇有挑战。但这样的努力值得：简明描述研究问题与设计，可以澄清研究者的思路，也有助于同事和顾问迅速了解研究。

\subsection{研究参与者}
\sourcepage{22}
选择研究参与者的样本时，需要作出两项主要决定（第 3 章）。首先要描述希望研究哪一类人，并规定界定参与者的\term{纳入标准（inclusion criteria）}和\term{排除标准（exclusion criteria）}。其次，要决定如何从该人群中一个可接触的子集招募适当数量的参与者。例如，关于咖啡因摄入与痴呆的病例对照研究，在取得 IRB 批准后，可以先查阅所在机构的电子健康记录，找出诊断为痴呆或认知功能障碍的患者，再联系他们及适当的对照，询问是否愿意参加。

\subsection{变量}
设计任何研究时，另一组重要决定是选择测量哪些\term{变量（variables）}（第 4 章）。例如，咖啡因摄入研究可能询问含不同咖啡因量的各种饮料和提神药片，并询问每份摄入量、摄入频率及时间。

分析性研究考察变量之间的关联，以预测结局并作出因果推断。考虑两个变量的关联时，发生较早，或从生物学上更可能起因果作用的变量，称为\term{预测变量（predictor variable）}；另一个称为\term{结局变量（outcome variable）}\footnote{预测变量有时称为自变量（independent variables），结局称为因变量（dependent variables），但这些名称的含义不那么直观。预测变量可以是暴露、危险因素或保护因素、治疗或干预，也可以是检验结果；全书会在适当语境中使用这些名称。}。多数观察性研究有许多预测变量，例如年龄、性别、种族／族裔、吸烟史和咖啡因摄入量，也有若干结局变量，例如认知功能与生活质量（quality of life）。

临床试验考察干预的效果。干预是一类特殊的预测变量，由研究者加以操控，例如给予咖啡因胶囊或相应的安慰剂。这种设计使她能利用\term{随机化（randomization）}尽量减少\term{混杂变量（confounding variables）}的影响。混杂变量是结局的其他预测因素，例如吸烟或受教育程度，它们也影响咖啡因摄入，因而可能使研究发现的解释发生混淆。混杂（confounding）尤其棘手，第 10 章将详细讨论。

\sourcepage{23}
\subsection{统计学问题}
研究者必须制订样本量估算、研究数据管理与分析计划。这通常需要提出\term{研究假设（research hypothesis）}（第 5 章），即把研究问题表述为可检验研究发现\term{统计学显著性（statistical significance）}的形式：
\begin{studyquote}
假设：在 60—74 岁成年人中，平均每天喝至少两杯咖啡（或摄入等量咖啡因）者的认知功能，优于摄入量较低者。
\end{studyquote}
研究假设还使研究者能够计算样本量：假如研究组之间确有结局差异，为以合理概率检出统计学显著差异，所需的参与者人数。这样的检出能力称为\term{检验效能（power）}（第 5 章）。纯描述性研究，例如“认知功能正常者中有多大比例每天喝咖啡”，不涉及统计学显著性检验，因此不需要假设；此时可以计算所需人数，使均值、比例或其他描述性统计量的\term{置信区间（confidence intervals）}足够窄。

\section{研究的运作：它如何发挥作用}
临床研究的目标，是利用研究发现增进对现实世界的理解（见图 1.1 下半部）。首先必须解释研究发现；这一过程取决于\term{内部效度（internal validity）}，即研究结果在多大程度上反映研究中实际发生的情况。随后，研究者从研究结论作出推断，增进对现实世界的理解。这一过程取决于\term{外部效度（external validity）}，或\term{可推广性（generalizability）}，即这些结论能够在多大程度上适用于研究之外的人与事。

\sourcepage{24}
\par
\begin{figure}[htbp]\centering\chonebasicflow
\caption{临床研究的基本结构。每一步中的选择与误差，都会影响研究完成后能够作出的推断。}\label{fig:01_01}\end{figure}\par
研究者在规划研究时，会反向进行这个过程，沿图 1.1 上半部从左到右，目标是在研究结束时使这些推断尽可能有效。她在研究计划中选择研究问题、参与者与测量，使外部效度提高，并使研究便于实施，以尽可能提高内部效度。下面先讨论设计，再讨论实施，最后讨论威胁这些推断效度的误差。

\subsection{设计研究}
请考虑下面这个关于咖啡因与认知功能的简单描述性问题：
\begin{studyquote}
70 岁及以上、认知功能正常的成年人，每天平均摄入多少咖啡因？
\end{studyquote}
这个问题不可能得到完全准确的答案，因为不可能研究由全体老年人构成的\term{目标人群（target population）}，更不用说咖啡因摄入量与认知功能的测量方法也不完美。因此，研究者只能退而求其次，回答一个与之相关、但能够作答的问题：
\begin{studyquote}
向在研究者所在机构就诊、年龄至少 70 岁、且电子病历中未记录痴呆或认知功能障碍诊断的患者邮寄问卷，得到的自报平均每日咖啡因摄入量是多少？
\end{studyquote}

\sourcepage{25}
图 1.2 展示了从研究问题转化为研究计划的过程。一个主要环节，是选择代表人群的参与者\term{样本（sample）}。由于研究整个目标人群几乎总会面临实际障碍，方案规定的研究群体通常只能是该人群的一个样本。利用研究者所在机构寻找潜在参与者，是使研究可行的折中，但它也有缺点：即使每个人都寄回问卷，所得的咖啡因摄入模式仍可能不同于一般人群。更重要的是，没有记录痴呆或认知功能障碍诊断，绝不意味着一个人的认知功能必然正常。
\par
\begin{figure}[htbp]\centering\chonedesignflow
\caption{研究设计中的选择与误差影响外部效度。如果预定样本与变量不能充分代表目标人群和所关注的现象，就会削弱通过推断增进对问题理解的能力。}\label{fig:01_02}\end{figure}\par

\sourcepage{26}
这一转化的另一个主要环节，是选择能够代表所关注\term{现象（phenomena）}的变量。研究计划规定的变量通常只是这些现象的替代指标。例如，用自报问卷评估过去一个月的咖啡因摄入量，收集信息既快又便宜。但它不可能完全准确，因为人们通常记不清上个月究竟摄入了多少咖啡因，而且那个月也未必反映他们平常的摄入模式。

因此，图 1.2 中研究问题与研究计划的每一项差别，都代表了为了提高可行性和效率而作的选择。然而，这种便利有其代价：这些选择可能使研究实际回答的问题与原先的问题有所不同，因而降低结论的相关性。

\subsection{实施研究}
研究计划的实施，影响实际研究忠实执行计划的程度。这里的问题是：如果样本抽取或测量的实际做法，与设计有很大差别，就可能对研究问题给出错误答案（图 1.3）。
\sourcepage{27}
\par
\begin{figure}[htbp]\centering\choneimplementationflow
\caption{研究实施中的选择与误差影响内部效度。如果实际样本与测量不能充分代表预定样本与变量，就会削弱从研究发现得出有效结论的能力。}\label{fig:01_03}\end{figure}\par
参与者的实际样本几乎总是不同于预定样本。例如，原计划研究门诊中所有符合条件、无痴呆或认知功能障碍的患者，但电子病历中的诊断不完整、邮寄地址错误以及拒绝参加等问题，都可能打乱计划。能够联系到且同意参加者的咖啡因摄入模式，可能不同于未联系到或不感兴趣者。除这些问题外，实际测量也可能偏离预定测量。如果问卷格式不清楚，参与者可能混淆而勾错方框，或无意中漏答一道题。

\sourcepage{28}
\subsection{因果推断}
许多研究旨在确立因果关联，也就是预测因素导致结局，以识别可能改善健康的临床或公共卫生干预。例如，如果摄入咖啡因会降低发生痴呆的风险，也许我们就应广泛推荐。试图确立\term{因果推断（causal inference）}的研究，会遇到一种特殊的效度问题。如果队列研究发现咖啡因摄入与痴呆相关，这代表因果关系，还是咖啡因摄入仅与另一种具有痴呆保护作用的混杂暴露——例如阅读晨报——有关？减少混杂及其他竞争性解释的可能性，是观察性研究设计中的主要挑战之一（第 10 章）。

\subsection{研究中的误差}
既然没有研究完全没有误差，目标就是让依据研究观察作出的各种推断，具有尽可能高的效度。有些错误推断可以在分析阶段处理，但更好的策略是着眼于设计与实施（图 1.4），从一开始就防止误差发生。

\sourcepage{29}
\par
\begin{figure}[p]\centering\chonefullflow
\caption{研究的运作。设计与实施中的选择和误差，会影响研究的内部效度与外部效度。}\label{fig:01_04}\end{figure}\par
妨碍研究推断的错误，主要分为\term{随机误差（random error）}和\term{系统误差（systematic error）}。区分两者很重要，因为减少它们的策略截然不同。

随机误差是偶然造成的错误，意思是找不到已知原因，也没有可预测的误差模式。随机误差可以向任一方向扭曲测量。譬如，在某医保计划数千名 70 岁及以上成员中，每天摄入咖啡因的真实比例恰为 40\%。如果从这个\term{可接触人群（accessible population）}中设计良好地抽取 100 人，样本可能恰好有 40 人每天摄入咖啡因。但更可能出现稍有不同的人数，例如 38、39、41 或 42 人。偶然因素有时也会产生差异较大的人数，例如 32 或 49 人。减少随机误差影响的方法很多，最简单的是增加样本量。\sourcepage{30}较大的样本可提高估计\term{精度（precision）}，也就是每次抽样所得比例接近 40\% 的程度，从而降低结果严重偏离的可能性。但增加样本量也会提高成本。幸而还有其他减少随机误差的方法，包括使用更好的测量工具（第 4 章）。

系统误差是\term{偏倚（bias）}导致的错误结果：某些变异来源，使研究发现朝一个方向偏离。图 1.2 中选择研究某个本地医保计划的患者，就是一个例子；该计划为避免——或鼓励！——过度编码所作的努力，可能影响痴呆与认知功能障碍诊断的记录。增加样本量对系统误差没有作用。提高估计\term{准确度（accuracy）}，也就是其接近真实值的程度，最佳方法是减小潜在偏倚的幅度；本书很大一部分内容都讨论这一点。此外，研究者还可以寻找补充信息，评价可能存在的偏倚有多重要，例如在另一场所再抽取一组样本。

上面两段的随机误差与系统误差示例，都属于\term{选择偏倚（selection bias）}的组成部分，会威胁从研究参与者推广到人群的推断。随机误差和系统误差也都可以造成\term{测量误差（measurement error）}，威胁从研究测量推广到所关注现象的推断。随机测量误差的例子是同一问卷多次施测时，回答有所变化。系统测量误差的例子是问卷存在问题，例如未询问能量饮料，导致咖啡因摄入量估计错误。第 3、4 章将介绍控制这些误差的策略。

如表 1.3 所概括，要尽可能回答好研究问题，就需要在研究设计、实施与解释中，尽量减小潜在误差。

\sourcepage{31}
\begin{center}\small
\captionof*{table}{\textbf{表 1.3\quad 研究设计与实施阶段控制随机误差及系统误差的方法}}
\begin{tabularx}{\linewidth}{L{37mm}Y}\toprule
\textbf{误差类型}&\textbf{解决方法}\\\midrule
随机误差&提高测量的精密度（第 4 章）；增加样本量（第 5、6 章）。\\
系统误差&提高测量的准确度与效度（第 4 章）；选择更好的设计（第 8—14 章）。\\
随机误差及系统误差&质量控制（第 18 章）。\\\bottomrule
\end{tabularx}\end{center}

\section{转化研究}
\term{转化研究（translational research）}是指将基础科学发现转化为临床实践，或将临床研究结果拓展为大规模公共卫生或社区项目的研究。其设计方法与其他临床研究相同，但参与的研究者往往更多，过程也更复杂，并需反复调整。成功的转化研究要求研究链条各环节的研究者合作，而且人人都应准备修改计划。在实验动物中看来很有希望的治疗，用于人体时产生不良反应；在学术门诊表现良好的筛查试验，在社区中使用繁琐或准确度下降——这些情况并不少见。合作越早开始，研究者越有可能提出重要问题、制订有价值的计划，并避免本可预见的挫折。

\subsection{从实验室研究转化到临床研究}
DNA 测序（DNA sequencing）、基因表达芯片（gene expression arrays）、分子影像（molecular imaging）、蛋白质组学（proteomics）、代谢组学（metabolomics）等许多工具，已从基础科学实验室进入临床研究。与普通临床研究相比，要成功转化实验室发现的检验和治疗方法，通常需要研究者本人有实验室经验，或有具备这些技能的合作者。“从实验台到病床”的研究\sourcepage{32}要求深入理解相关基础科学。许多临床研究者以为这些知识很容易掌握，正如许多实验室科学家以为临床研究不需要特殊训练和专业能力。然而，基础研究与临床研究的技能几乎不重叠，要提出好的转化研究问题，需要多个领域的专业能力与密切合作。

例如，假设已发现一种调节小鼠昼夜节律的基因，一位擅长睡眠障碍的临床研究者，希望研究人类同源基因的变异是否影响睡眠模式。要就该基因的潜在作用提出有意义的问题，她需要合作者熟悉该基因及其蛋白产物的生物学，也熟悉各种基因分型方法的优势与局限。再如，一位实验室研究者在乳腺癌患者的淋巴结中发现了独特的基因表达模式。若要研究该模式是否具有预后价值，她应与了解重测信度、抽样、盲法，以及疾病先验概率如何影响发现的适用性和临床用途的研究者合作。最后，一个想检验新药的团队，可能需要熟悉分子生物学、药代动力学与药效动力学、I 期和 II 期临床试验，以及相关医学领域诊疗模式的科学家。

\subsection{从临床研究转化到人群研究}
设计研究，将临床发现应用于更大、更多样的人群，也可能需要额外的专业能力，例如识别高危人群或医疗服务不足的人群、理解筛查与诊断的区别，以及知道如何在医疗服务系统中推行改变。在实践中，这类研究通常需要接触大规模参与者群体，例如某医保计划的参保者或全县居民。规划时，获得系主任、附属医院医疗团队负责人、\sourcepage{33}管理式医疗组织负责人或地方卫生部门主管的支持与建议，可能很有帮助。

一些研究者设计此类转化研究时会走捷径。例如，为拓展本专科门诊的研究结果，他们计划从其他大学医疗中心的相似门诊招募参与者，却不让社区医师参与。这有点像把阿里斯托芬的作品译为现代希腊语——对于英语读者，仍然没有多少用处。在更大、更多样的人群中检验研究发现，常常需要与社区成员合作，并调整方法，以适应非学术机构，见第 15 章。

\section{设计研究}
\subsection{研究计划}
制订研究计划，首先用一句话表述研究问题，明确主要预测变量、结局变量和目标人群。随后依次形成三种版本，每一种都比前一种更长、更详细。
\begin{itemize}
\item \textbf{研究提纲（study outline）}（表 1.1 与附录 1A）。这是用一页纸概括的设计，作为格式统一的清单，提醒研究者处理所有组成部分。它按有序的逻辑展开，有助于澄清研究者对主题的思考。
\item \textbf{研究方案。}在提纲基础上扩展，通常为 5—15 页，用于规划研究，申请 IRB 批准及科研资助。全书会逐项讨论其组成部分，第 20 章加以汇总。
\item \textbf{操作手册（operations manual）。}招募参与者并收集数据的研究，还要编写全面的手册，纳入具体操作说明、问卷、数据收集表和其他材料。即使使用事先收集的资料，例如电子健康记录，也必须准确规定怎样使用。手册使研究\sourcepage{34}以统一、标准化的方式实施，并有良好的质量控制（quality control）（第 18 章）。
\end{itemize}
应在早期把研究问题和提纲写下来。将想法付诸文字，有助于把模糊构想变成具体计划，也为向同事与顾问征求意见提供具体基础。这并不容易——想法说起来比写下来容易——但回报是更快启动和更好的项目。有些初学者急于开始研究，不愿花时间写明详细计划，这可以理解，却常会导致项目构思不周和后续延误。

附录 1A 是一个研究提纲示例。这样的提纲更多涉及研究的结构（表 1.1），而非运作（图 1.4），所以研究者必须提醒自己，仔细考虑怎样尽量减少从研究发现作出推断时可能出现的误差。根据拟采用的参与者获取与测量计划，以及可能发生的实施问题，明确比较研究实际可能回答的问题与原研究问题有何差别，就可以发现研究的长处与问题。

有了提纲，并明确预定推断之后，研究者即可继续拟定方案细节。这是反复调整的过程：征求同事意见，起草具体招募和测量方法，考虑科学与伦理上的适当性，必要时修改研究问题与提纲，预试具体招募及测量方法，再作修改，再听取意见，如此继续。

\subsection{取舍}
遗憾的是，误差是所有研究不可避免的一部分。主要问题在于，误差是否大到使研究结论不可靠。设计研究时，研究者的处境很像一名为新合同谈判的工会代表。她开始时列出愿望清单：缩短工时、提高工资、改善福利等。随后必须作出让步，坚持最重要的内容，\sourcepage{35}放弃非必需或不现实的要求。谈判结束，还有最后一个至关重要的步骤：审视自己能争取到的最佳合同，决定是否值得签署。

研究者把研究问题转化成计划，并考虑潜在实施问题时，也必须作类似让步。一边是内部与外部效度，另一边是可行性，总要有所折中。但是，研究计划形成后，研究者应重新审视：它是否充分回答了研究问题？实施时能否把误差控制在可接受水平？最后决定按这个设计是否值得开展研究，相当于工会谈判者决定是否接受合同。遗憾的是，这一步有时被省略；有时答案是否定的，就需要重新开始。不过不必灰心！好的科学家之所以出色，并非他们的想法个个都好，而是他们能迅速放下行不通的想法，转向更好的想法。

\section{小结}
\begin{enumerate}
\item 研究的结构是构成研究计划的一组具体要素：研究问题及其意义，以及设计、研究参与者和测量方法。挑战是设计这些要素，在可行性与得出有效结论的需要之间作出权衡。
\item 研究的运作，是研究怎样发挥作用。研究者根据发现，对可接触人群中可测量的现象得出结论（内部效度），再将结论应用于其他情境、人物与事件（外部效度或可推广性）。挑战是设计并实施一个计划，充分控制威胁效度的两大因素——随机误差（偶然性）与系统误差（偏倚），并增进有关研究问题的现有知识。
\item 设计研究时，参照图 1.4，思考研究问题（她希望回答什么）、\sourcepage{36}研究计划（设计要回答什么），以及实际研究（考虑预期实施误差后，实际上会回答什么）之间的关系，可能很有帮助。
\item 制订计划的一种好方法，是先用一句话写明主要变量和目标人群，再扩展成一页提纲，以标准顺序列出研究要素。此后再扩展为研究方案与操作手册。
\item 在各种取舍中，以及判断项目整体是否可行时，都需要研究者的良好判断和同事的建议。
\end{enumerate}

\chapterappendix{1A}{研究提纲示例}
\sourcepage{37}
下面是一项研究项目的简要计划，由 Michael Jung, MD, MBA 开展。他在 UCSF 担任麻醉科住院医师期间，在 Jina Sinskey, MD 指导下启动这项研究。多数初学研究者觉得观察性研究更容易开展，但在这一情形中，一项规模与范围适中的随机临床试验是可行的。研究结果发表于一份影响较大的期刊（1）。

\section{儿童麻醉围术期虚拟现实干预}
\subsection{背景与意义}
儿童围术期焦虑，是麻醉与手术对儿童心理影响中常见且重要的一面。既往研究表明，接受手术及全身麻醉的儿童中，麻醉诱导时高度焦虑的发生率可高达 50\%（2）。视听分散注意力的方法，包括虚拟现实（virtual reality，VR），有望成为安全、无创、非药物性的减轻焦虑手段。

\subsection{具体目标}
开展一项随机对照试验（randomized controlled trial），考察 VR 头戴设备提供的沉浸式视听分心干预，对接受择期手术及全身麻醉患儿围术期焦虑的影响。

\subsection{方法}
\subsubsection{设计概述}
在这项前瞻性、随机、对照、平行组研究中，计划接受择期手术及全身麻醉的儿童，将被随机分配至 VR 组或对照组（不使用 VR）。VR 组患儿在手术室全身麻醉诱导期间，接受 VR 视听分心干预。使用已验证的改良耶鲁术前焦虑量表（Modified Yale Preoperative Anxiety Scale，mYPAS）评价焦虑。

\sourcepage{38}
\subsubsection{研究参与者}
本项目的目标人群，是接受择期手术及全身麻醉的 5—12 岁儿童；可接触人群是 UCSF Benioff 儿童医院就诊的同类儿童。

\subsubsection{测量}
预测变量是干预，即使用或不使用 VR 头戴设备。主要结局是儿童围术期焦虑（以 mYPAS 评分测量），在三个时间点评价：术前等候区、进入手术室时，以及全身麻醉诱导期间。

\subsubsection{随机化与盲法}
采用计算机生成的随机分配表，将其纳入已编程的研究应用程序 Redcap，并对研究招募人员隐藏。为进行随机分配，录入带有时间戳的记录，包括病历号及完成知情同意的信息。鉴于 VR 头戴设备的性质——患者要么戴着，要么没有戴着——不采用盲法。

\subsubsection{数据分析}
预测变量是二分类变量：使用或不使用 VR 头戴设备。结局变量为麻醉诱导时的 mYPAS 评分，是取值 0—100 的连续变量，将用 $t$ 检验比较两组。分类结局使用卡方分析，显著性水平设为 0.05。

\subsubsection{样本量估算}
根据一项关于 5—12 岁儿童术前焦虑的既往研究（3），对照组 mYPAS 评分的均值 $\pm$ 标准差为 $30.1\pm8.4$。每组招募 31 名参与者，可以 80\% 的检验效能（$\beta=0.20$），检出两组在诱导时 mYPAS 评分 20\% 的差异。

\section*{附录 1A 参考文献}
\sourcepage{39}
\begingroup\small\begin{enumerate}[leftmargin=2em,itemsep=3pt]
\item Jung MJ, Libaw JS, Ma K, Whitlock EL, Feiner JR, Sinskey JL. Pediatric distraction on induction of anesthesia with virtual reality and perioperative anxiolysis: a randomized controlled trial. \textit{Anesth Analg}. 2021;132(3):798-806.
\item Kain ZN, Mayes LC, Caldwell-Andrews AA, Karas DE, Mcclain BC. Preoperative anxiety, postoperative pain, and behavioral recovery in young children undergoing surgery. \textit{Pediatrics}. 2006;118(2):651-658.
\item Moura LA, Dias IM, Pereira LV. Prevalence and factors associated with preoperative anxiety in children aged 5-12 years. \textit{Rev Lat Am Enfermagem}. 2016;24:e2708.
\end{enumerate}\endgroup

\chapterappendix{1B}{第 1 章习题：起步——临床研究的结构与运作}
\sourcepage{40}
\begin{enumerate}
\item 早期少量配方奶（Early Limited Formula，ELF）研究在加利福尼亚州两所学术医疗中心开展，目的是促进出生后 36 小时内体重下降 $\geq5\%$ 的新生儿母乳喂养。该随机临床试验中，ELF 干预是教父母在每次母乳喂养后，用注射器喂 10 mL 配方奶，直到成熟乳开始分泌；对照组父母学习安抚婴儿的技巧。3 个月时，向不知分组的访谈员报告纯母乳喂养的母亲比例，ELF 组为 79\%，对照组为 42\%（$P=0.02$）（1）。

对以下各项陈述，请指出：（a）这是有关内部效度还是外部效度的推断；（b）你认为该推断是否有效；（c）可能无效的理由。
\begin{enumerate}
\item 对本研究中的妇女，提供 ELF 提高了 3 个月时的母乳喂养率。
\item 对出生于波士顿一家社区医院、出生后 36 小时内体重下降 $\geq5\%$ 的婴儿，提供 ELF 很可能提高 6 月龄时的母乳喂养率。
\item 根据本研究结果，开展国际行动向多数新生儿提供配方奶，很可能促进成功母乳喂养，并改善新生儿及其母亲的健康。
\end{enumerate}
\item 请根据以下已发表研究的摘要，各用一句话说明研究设计与研究问题，并纳入主要预测变量、结局变量及预定样本。
\begin{enumerate}
\item 北卡罗来纳州温斯顿-塞勒姆的研究者，随机抽取当地 2228 名高中生，询问过去两周观看电视摔角节目的频率；6 个月后，再询问同一批学生在校及约会时打架的情况。学生报告 6 个月前每多看一集摔角节目，报告曾与约会对象打架的经调整比值（odds）就增加 14\%（2）。
\item 为评估母乳喂养量是否保护妇女免于卵巢癌，研究者调查了 493 名新诊断卵巢癌的中国妇女，以及 472 名其他住院妇女；所有人都曾母乳喂养至少一个孩子。结果发现，母乳喂养累计月数与卵巢癌风险降低存在剂量－反应关系。\sourcepage{41}例如，与母乳喂养不足 10 个月者相比，至少喂养 31 个月者的比值比（odds ratio）为 0.09（95\% 置信区间：0.04—0.19）（3）。
\item 为了解不育男性中膳食饱和脂肪摄入与精子浓度降低的关联，是否也适用于一般人群，丹麦研究者在年轻男性接受征兵体检时，向同意参与者收集精液样本和食物频率问卷。结果发现，自报膳食饱和脂肪摄入量与精子浓度降低存在显著的剂量－反应关系。例如，与摄入量最低四分位组相比，最高四分位组的精子浓度低 41\%（95\% 置信区间：4\%、64\%）（4）。
\item 通过封堵左心耳，或可预防心房颤动中的栓塞性卒中。研究者纳入因其他适应证接受心脏手术的心房颤动患者（5），在手术中对随机选出的一半患者封堵左心耳，另一半不封堵。封堵组 2379 人中，114 人（4.8\%）发生卒中或全身性栓塞；未封堵组 2391 人中，168 人（7.0\%）发生上述结局（$P<0.001$）。
\end{enumerate}
\end{enumerate}

\begin{references}
\item Flaherman VJ, Aby J, Burgos AE, Lee KA, Cabana MD, Newman TB. Effect of early limited formula on duration and exclusivity of breastfeeding in at-risk infants: an RCT. \textit{Pediatrics}. 2013;131(6):1059-1065.
\item DuRant RH, Champion H, Wolfson M. The relationship between watching professional wrestling on television and engaging in date fighting among high school students. \textit{Pediatrics}. 2006;118:e265-e272.
\item Su D, Pasalich M, Lee AH, Binns CW. Ovarian cancer risk is reduced by prolonged lactation: a case-control study in southern China. \textit{Am J Clin Nutr}. 2013;97:354-359.
\item Jensen TK, Heitmann BL, Jensen MB, et al. High dietary intake of saturated fat is associated with reduced semen quality among 701 young Danish men from the general population. \textit{Am J Clin Nutr}. 2013;97:411-418.
\item Whitlock RP, Belley-Cote EP, Paparella D, et al. Left atrial appendage occlusion during cardiac surgery to prevent stroke. \textit{New Engl J Med}. 2021;384:2081-2091.
\end{references}
\sourcepage{42}
